% Options for packages loaded elsewhere
\PassOptionsToPackage{unicode}{hyperref}
\PassOptionsToPackage{hyphens}{url}
\PassOptionsToPackage{dvipsnames,svgnames,x11names}{xcolor}
%
\documentclass[
  11pt,
  a4paper,
]{ltjsarticle}
\usepackage{amsmath,amssymb}
\usepackage{iftex}
\ifPDFTeX
  \usepackage[T1]{fontenc}
  \usepackage[utf8]{inputenc}
  \usepackage{textcomp} % provide euro and other symbols
\else % if luatex or xetex
  \usepackage{unicode-math} % this also loads fontspec
  \defaultfontfeatures{Scale=MatchLowercase}
  \defaultfontfeatures[\rmfamily]{Ligatures=TeX,Scale=1}
\fi
\usepackage{lmodern}
\ifPDFTeX\else
  % xetex/luatex font selection
\fi
% Use upquote if available, for straight quotes in verbatim environments
\IfFileExists{upquote.sty}{\usepackage{upquote}}{}
\IfFileExists{microtype.sty}{% use microtype if available
  \usepackage[]{microtype}
  \UseMicrotypeSet[protrusion]{basicmath} % disable protrusion for tt fonts
}{}
\makeatletter
\@ifundefined{KOMAClassName}{% if non-KOMA class
  \IfFileExists{parskip.sty}{%
    \usepackage{parskip}
  }{% else
    \setlength{\parindent}{0pt}
    \setlength{\parskip}{6pt plus 2pt minus 1pt}}
}{% if KOMA class
  \KOMAoptions{parskip=half}}
\makeatother
\usepackage{xcolor}
\usepackage[margin=24mm]{geometry}
\setlength{\emergencystretch}{3em} % prevent overfull lines
\providecommand{\tightlist}{%
  \setlength{\itemsep}{0pt}\setlength{\parskip}{0pt}}
\setcounter{secnumdepth}{5}
\ifLuaTeX
  \usepackage{selnolig}  % disable illegal ligatures
\fi
\IfFileExists{bookmark.sty}{\usepackage{bookmark}}{\usepackage{hyperref}}
\IfFileExists{xurl.sty}{\usepackage{xurl}}{} % add URL line breaks if available
\urlstyle{same}
\hypersetup{
  colorlinks=true,
  linkcolor={blue},
  filecolor={Maroon},
  citecolor={Blue},
  urlcolor={blue},
  pdfcreator={LaTeX via pandoc}}

\author{}
\date{}

\begin{document}

\hypertarget{ux6f14ux7fd2-07-ux884cux5217ux5206ux89e3}{%
\section{演習 07 ---
行列分解}\label{ux6f14ux7fd2-07-ux884cux5217ux5206ux89e3}}

\hypertarget{lu}{%
\subsection{1. LU}\label{lu}}

\[
A=\begin{bmatrix}2&1\\4&3\end{bmatrix}
\]

を \texttt{LU} 分解してください。

\hypertarget{qr}{%
\subsection{2. QR}\label{qr}}

\[
A=\begin{bmatrix}1&1\\1&0\\0&1\end{bmatrix}
\]

の列に Gram--Schmidt を適用して reduced QR 分解を求めてください。

\hypertarget{svd-ux306eux57faux672cux5f0f}{%
\subsection{3. SVD の基本式}\label{svd-ux306eux57faux672cux5f0f}}

\texttt{A=UΣV\^{}T} から

\[
A^TA=VΣ^TΣV^T
\]

を導出してください。

\hypertarget{ux30e9ux30f3ux30af}{%
\subsection{4. ランク}\label{ux30e9ux30f3ux30af}}

SVD から \texttt{rank(A)}
が非零特異値の個数であることを説明してください。

\hypertarget{ux5e7eux4f55}{%
\subsection{5. 幾何}\label{ux5e7eux4f55}}

\texttt{A}
が単位球を楕円体へ写すとき、右特異ベクトル、左特異ベクトル、特異値がそれぞれ何を意味するか説明してください。

\hypertarget{ux6bd4ux8f03}{%
\subsection{6. 比較}\label{ux6bd4ux8f03}}

固有値分解と SVD
の相違点を「対象行列」「基底」「値の符号」「幾何学」の4観点から整理してください。

\end{document}
